\section{An integral null class with a uniform chamber}
\label{sec:lattice-chamber}

The curve-existence result of \Cref{thm:family-curves} applies to every
nonzero integral null class.  We now choose one for which a positive curve
configuration cannot split.  The choice depends only on the original
period plane, so the resulting restrictions on curve classes persist
under later deformations of the definite pair that preserve its two
periods.
The arithmetic step is to find a primitive null class whose projection
direction avoids a finite set of obstructions.  Adjunction then relates
this arithmetic choice to curve classes.  Taming and the resulting class alternatives rule out splitting or
multiplicity in a positive configuration in the chosen class; positivity
of intersections makes distinct curves in that class disjoint.

We use the free lattice
\[
  \Lambda=H^2(X;\Z)/\mathrm{torsion},\qquad
  \Lambda_{\Q}=\Lambda\otimes_{\Z}\Q,\qquad
  \Lambda_{\R}=\Lambda\otimes_{\Z}\R.
\]
A curve class will mean the image in \(\Lambda\) of its integral
Poincar\'e dual, with the complex orientation of the curve.  Equalities
of integral classes in this section are therefore equalities modulo
torsion.  This convention does not affect intersections or integrals of
closed real forms.  The original positive period plane is denoted by
\(P\), and \(d^+\) denotes the orthogonal projection of
\(d\in\Lambda_{\R}\) to \(P\).  By \Cref{prop:topology},
\[
  \Lambda\simeq 3U\oplus 2E_8(-1)
  \quad\hbox{or}\quad
  \Lambda\simeq 3U,
\]
where \(U\) is the hyperbolic plane.  In particular, \(P^\perp\) is
negative definite.  No rationality of \(P\) is assumed.

\begin{proposition}[Choice of chamber class]
\label{prop:chamber-class}
There is a primitive class \(F\in\Lambda\) with \(F^2=0\) such that the
following implication holds for every \(d\in\Lambda\) and every real
number \(c>0\):
\begin{equation}
\label{eq:lattice-chamber-implication}
  d^2\geq-2,\qquad d^+=cF^+,\qquad d\cdot F\leq0
  \quad\Longrightarrow\quad
  c\in\Z_{>0},\qquad d\cdot F=0.
\end{equation}
Moreover, \(F^+\ne0\).
\end{proposition}

\begin{proof}
We first construct a rational null subspace with enough integral
directions, and then exclude finitely many of those directions.

\smallskip
\noindent\emph{A positive integral vector and a null lattice.}
Set
\[
  W_{\Q}=\Lambda_{\Q}\cap P^\perp,\qquad
  W_{\R}=\operatorname{span}_{\R}W_{\Q}.
\]
The subspace \(W_{\R}\) is rational and negative definite.  Hence its
orthogonal complement is rational and contains \(P\).  Positive square
is an open condition, so density of rational vectors in this complement
gives a positive rational vector in \(W_{\R}^\perp\).  Clearing
denominators and dividing out the common integral divisor produces a
primitive vector \(l\in\Lambda\) satisfying
\begin{equation}
\label{eq:lattice-positive-vector}
  l^2>0,\qquad l\cdot W_{\Q}=0.
\end{equation}

Write \(l^2=2n\), where \(n\in\Z_{>0}\).  We use the following form of
Eichler's criterion: in an even lattice containing two orthogonal
hyperbolic planes, the orbit of a primitive vector under the stable
orthogonal group is determined by its square and its normalized class
in the discriminant group
\cite[Proposition~3.3(i)]{GritsenkoHulekSankaran2009}.
Here the lattice is unimodular.  Thus its discriminant group is zero,
and every primitive vector has divisibility one.  The criterion sends
\(l\) to \(e_1+n f_1\), where
\[
  e_i^2=f_i^2=0,\qquad e_i\cdot f_i=1
\]
are standard bases of the three mutually orthogonal copies of \(U\).
In these coordinates the lattice generated by \(f_1,e_2,e_3\) is
primitive and totally null.  Pulling it back gives a rational
three-dimensional totally null subspace \(L_{\R}\) and its lattice
\(L=\Lambda\cap L_{\R}\), with an element \(f_0\in L\) such that
\begin{equation}
\label{eq:lattice-affine-slice}
  l\cdot f_0=1.
\end{equation}
Consequently
\[
  \mathcal F=\{F\in L:l\cdot F=1\}
\]
is an affine lattice of rank two.  Every \(F\in\mathcal F\) is primitive
in \(\Lambda\), since a nontrivial integral divisor of \(F\) would
divide \(l\cdot F=1\).  It is also nonzero and null.  Its projection to
\(P\) is nonzero by negative definiteness of \(P^\perp\).

\smallskip
\noindent\emph{Only finitely many cosets can obstruct a direction.}
The set
\[
  K=\{r\in P^\perp:r^2\geq-2\}
\]
is compact.  Suppose temporarily that \(F\in\mathcal F\), \(d\in\Lambda\),
and \(c>0\) satisfy the three conditions on the left side of
\Cref{eq:lattice-chamber-implication}.  Then
\begin{equation}
\label{eq:lattice-bounded-residual}
  r=d-cF\in P^\perp,\qquad
  r^2=d^2-2c(d\cdot F)\geq-2,
\end{equation}
so \(r\in K\).

Let \(q:\Lambda_{\R}\to\Lambda_{\R}/L_{\R}\) be the quotient map.
Because \(L_{\R}\) is rational and \(L\) is saturated, a basis of \(L\)
extends to an integral basis of \(\Lambda\); the remaining basis
vectors project to a lattice basis in the real quotient.  Thus
\(q(\Lambda)\) is discrete and \(\ker(q|_\Lambda)=L\).
Since \(q(d)=q(r)\in q(K)\), only finitely many cosets \(d+L\) can occur
in \Cref{eq:lattice-bounded-residual}.  This finite set is independent
of \(F\).

\smallskip
\noindent\emph{Avoiding irrational points in a torus.}
Orthogonal projection restricts to an isomorphism
\[
  \pi:L_{\R}\longrightarrow P:
\]
its kernel is null and lies in the negative-definite space \(P^\perp\),
and both spaces have dimension three.  Hence \(\pi(L)\) is a lattice
in \(P\).  Each of the finitely many cosets just found determines one
point of
\[
  \mathbb T=P/\pi(L),
\]
because adding an element of \(L\) changes the projection by an
element of \(\pi(L)\).

For a primitive \(F\in L\), let
\[
  \mathbb S_F=\{t\pi(F)+\pi(L):t\in\R\}\subset\mathbb T.
\]
Two such circles with different directions intersect only in torsion
points.  Indeed, if
\(t\pi(F)-s\pi(F')\in\pi(L)\) and \(F,F'\) are not collinear, choose
an integral linear functional on \(L\) which vanishes on \(F\) and is
nonzero on \(F'\).  Applying it to
\(tF-sF'\in L\) shows that \(s\in\Q\), so the common point is torsion.
In particular, any non-torsion point of \(\mathbb T\) lies on at most
one primitive direction circle.

The elements of \(\mathcal F\) have distinct directions: two collinear
elements with \(l\)-pairing one are equal.  Since \(\mathcal F\) is
infinite, we can choose \(F\in\mathcal F\) whose circle contains none
of the non-torsion points among the finitely many possible coset
points.

For this \(F\), every \(d,c\) satisfying the hypotheses of
\Cref{eq:lattice-chamber-implication} gives a torsion point
\[
  d^++\pi(L)=c\pi(F)+\pi(L).
\]
There is therefore an integer \(k>0\) with \(kcF\in L\).  Pairing with
\(l\) shows that \(kc\in\Z\), so \(c\in\Q\).  It follows that
\(r=d-cF\in W_{\Q}\), and \Cref{eq:lattice-positive-vector} now gives
\begin{equation}
\label{eq:lattice-integral-coefficient}
  c=l\cdot d\in\Z_{>0}.
\end{equation}
If \(m=d\cdot F<0\), then \(m\) is a negative integer and
\[
  r^2=d^2-2cm\geq-2+2c\geq0.
\]
Negative definiteness of \(P^\perp\) forces \(r=0\), which would imply
\(m=0\).  This contradiction proves \(d\cdot F=0\), and completes the
proof.
\end{proof}

Fix \(F\) as in \Cref{prop:chamber-class}.  Make a constant orthogonal
change of the coefficient basis of the original triple and denote the
result by \((A,C,B)\), choosing
\begin{equation}
\label{eq:lattice-adapted-periods}
  [B]=\frac{F^+}{\lVert F^+\rVert},\qquad
  [A]\cdot F=[C]\cdot F=0,\qquad
  \beta=[B]\cdot F=\lVert F^+\rVert>0.
\end{equation}
Here the norm is the positive intersection norm on \(P\).  The three
classes remain orthonormal, and
\[
  V=[A]^\perp\cap[C]^\perp
\]
has signature \((1,b^-)\).  The following conclusion is deliberately
stated for arbitrary later definite pairs with these two periods.

\begin{corollary}[Curve-class alternatives]
\label{cor:curve-signs}
Let \(\widetilde A,\widetilde C\) be any smooth closed definite pair on
\(X\) with
\[
  [\widetilde A]=[A],\qquad [\widetilde C]=[C],
\]
and let \(J\) be either of its associated almost complex structures.
For every nonconstant irreducible closed \(J\)-holomorphic curve whose
class \(d\in\Lambda\) is nonzero, there is a real number \(c\ne0\) with
\(d^+=cF^+\).  Exactly one of the following alternatives holds:
\begin{enumerate}[label=\textup{(\roman*)}]
\item \(d\cdot F\ne0\), and
  \(\sgn(d\cdot F)=\sgn(c)=\sgn([B]\cdot d)\);
\item \(d=cF\) in \(\Lambda\), with \(c\in\Z\setminus\{0\}\).
\end{enumerate}
If \(J\) admits a closed taming form, every nonconstant irreducible closed
\(J\)-holomorphic curve has nonzero real class, so the nonzero-class
hypothesis is automatic.
\end{corollary}

\begin{proof}
Represent the irreducible image by a somewhere-injective
\(J\)-holomorphic map \(u:\Sigma\to X\), with its covering
multiplicity factored out.  The standard adjunction and singularity
theorems for simple curves in real dimension four give
\begin{equation}
\label{eq:lattice-adjunction}
  d^2=c_1(TX,J)\cdot d+2g(\Sigma)-2+2\delta(u),
  \qquad \delta(u)\in\Z_{\geq0},
\end{equation}
where \(\delta(u)\) is the weighted singularity defect and is zero
precisely for an embedding; see
\cite[Theorem~2.8]{Wendl2020}, using the numbering of the cited
prepublication version throughout.  Positivity of intersections of
distinct simple images is
\cite[Theorem~2.3 and Corollary~2.4]{Wendl2020}.
The critical points are isolated, and a simple curve is locally
injective, including at critical points
\cite[Propositions~B.26 and~B.41]{Wendl2020}.
These facts prevent accumulation of distinct double-point pairs at the
diagonal of the product of the source with itself.  Local intersection
isolation, unique continuation, and simplicity prevent accumulation off
the diagonal.
Compactness therefore makes the critical and noninjective sets finite.
The definite pair trivializes the canonical bundle by
\Cref{lem:definite-pair}, so \(c_1(TX,J)=0\) for either sign of \(J\).
In particular, \(d^2\geq-2\).

Both members of the pair vanish on complex lines.  Integrating their
pullbacks gives \([A]\cdot d=[C]\cdot d=0\); hence
\[
  d^+=cF^+,\qquad [B]\cdot d=\beta c
\]
for a real number \(c\).
If \(c=0\), the nonzero integral class \(d\) lies in \(P^\perp\).
Evenness, negative definiteness, and \(d^2\geq-2\) force \(d^2=-2\).
When \(b_1=0\), this contradicts \Cref{prop:root-exclusion}.
When \(b_1=4\), \Cref{eq:lattice-adjunction} gives
\(g(\Sigma)=\delta(u)=0\), so \(d\) is the class of an embedded sphere.
This contradicts the vanishing of spherical real classes in
\Cref{prop:cup-product}.  Thus \(c\ne0\).

Suppose first that \(c>0\).  By \Cref{prop:chamber-class},
\(d\cdot F<0\) is impossible.  If \(d\cdot F>0\), alternative
\textup{(i)} holds.  If \(d\cdot F=0\), the same proposition gives
\(c\in\Z_{>0}\), so \(r=d-cF\) is integral and lies in \(P^\perp\).
Moreover \(r^2=d^2\geq-2\).  Either \(r=0\), or negative definiteness
and evenness give \(r^2=d^2=-2\).  The latter is excluded by
\Cref{prop:root-exclusion} when \(b_1=0\); when \(b_1=4\), it would
make the original curve \(u\) an embedded sphere by
\Cref{eq:lattice-adjunction}, again contradicting
\Cref{prop:cup-product}.  Therefore \(d=cF\).

If \(c<0\), apply the same lattice argument to \(-d\), whose square
is still at least \(-2\).  This use of
\Cref{prop:chamber-class} is purely arithmetic and does not require
a \(J\)-holomorphic representative of \(-d\).  It gives
\(d\cdot F<0\) or \(d=cF\) with \(c\in\Z_{<0}\).
Indeed, a nonzero residual in \(P^\perp\) would again have square
\(d^2=-2\).  In the \(b_1=4\) case, adjunction is applied to the
original curve in class \(d\), making it an embedded sphere and giving
the same contradiction to \Cref{prop:cup-product}.
The two alternatives are disjoint because \(F^2=0\).

Finally, the integral of a closed taming form over any nonconstant
curve is strictly positive.  Its real homology class, and hence its
image in \(\Lambda\), cannot vanish.
\end{proof}

\Needspace{12\baselineskip}
\begin{proposition}[No splitting in the fiber class]
\label{prop:nonsplitting}
Let \((\widetilde A,\widetilde C,\widetilde B)\) be any smooth closed
positive triple with the same three cohomology classes as \((A,C,B)\),
and let \(J\) be the structure associated to its first two members
and tamed by \(\widetilde B\).  Then:
\begin{enumerate}[label=\textup{(\roman*)}]
\item Every positive \(J\)-holomorphic curve configuration whose
  total class is \(F\) in \(\Lambda\) consists of one reduced
  irreducible component of class \(F\).
\item Through every point of \(X\) there is such a component, and
  distinct components of class \(F\) are disjoint.
\item Every simple parametrization of such a component is either an
  embedded torus or a rational curve with total adjunction defect
  one.  If the rational parametrization is immersed, its image has
  exactly one transverse positive double point and no other
  singularity.
\end{enumerate}
These statements in particular hold after any exact perturbation of
the triple that preserves positivity.
\end{proposition}

\begin{proof}
Write a positive configuration as
\[
  F=\sum_{i=1}^r n_i d_i\quad\hbox{in }\Lambda,
  \qquad n_i\in\Z_{>0},
\]
where the \(d_i\) are the classes of its distinct irreducible images.
Taming excludes zero real component classes.  In the notation of
\Cref{cor:curve-signs},
\[
  0<\int_{d_i}\widetilde B=[B]\cdot d_i=\beta c_i,
\]
so \(c_i>0\).  The corollary gives \(m_i=d_i\cdot F\geq0\), whereas
\[
  0=F^2=\sum_{i=1}^r n_i m_i.
\]
Every \(m_i\) therefore vanishes.  Again by the corollary,
\(d_i=c_iF\) with \(c_i\in\Z_{>0}\), and comparison of the total
classes yields
\begin{equation}
\label{eq:lattice-no-splitting}
  \sum_{i=1}^r n_i c_i=1.
\end{equation}
There is exactly one summand, with \(n_1=c_1=1\).  This proves
\textup{(i)}, including the exclusion of a nontrivial covering
multiplicity.

Apply \Cref{thm:family-curves} to the present triple.  Its period
plane is still \(P\), and \Cref{eq:lattice-adapted-periods} forces the
direction of the configuration to be precisely the one whose
anti-invariant plane is
\(\operatorname{span}(\widetilde A,\widetilde C)\) and whose tamer is
\(\widetilde B\).  Thus the theorem gives a positive \(J\)-holomorphic
configuration in \(F\) through any prescribed point.  Part
\textup{(i)} makes it a single reduced component.  Two distinct
components of class \(F\) cannot meet: their intersection number is
\(F^2=0\), while any intersection of distinct simple images has
positive local intersection number.  This proves \textup{(ii)}.

For the remaining assertion, adjunction becomes
\[
  0=F^2=2g(\Sigma)-2+2\delta(u),
  \qquad\text{or equivalently}\qquad g(\Sigma)+\delta(u)=1.
\]
If \(g(\Sigma)=1\), the defect is zero and the curve is embedded.
If \(g(\Sigma)=0\), the total defect is one; critical points of the
parametrization have not yet been excluded.  When the map is
immersed, its singular defect is the sum of positive intersection
multiplicities between distinct local branches.  Total defect one
then forces exactly one pair of branches meeting with multiplicity
one, which is a transverse positive double point.  A tangency has
larger multiplicity, and a point with three or more branches
contributes at least three.  There can be no additional singularity.
\end{proof}
